\section{Data and design}\label{sec:data}

\paragraph{Patents and lapses.} The data are PATSTAT Global \citep[autumn 2025 edition;][]{PATSTAT}. The sample is
granted invention patents filed in 2000--2005 at ten national offices, followed for up to twenty years (to 2019 in Korea and 2022 in Russia, where the
lapse records end). A patent
lapses if the INPADOC legal-event table records its first cessation within 19.5 years of filing; a later cessation
marks expiry at full term.\footnote{In Brazil, where the term ran to at least ten years after grant, the cutoff is the
later of 19.5 years after filing and 9.5 years after grant, and cessations recorded before grant (unpaid annuities on
pending applications) are not counted.} I chose the offices by scanning every authority in the legal-event table for the share of
grants with a recorded cessation and the continuity of records over time (Appendix~\ref{app:data}). The
emerging-market offices are those of China, Korea, Russia, Brazil and the Czech Republic; the advanced-economy
offices are those of the United States, Japan, Germany, the United Kingdom and France. Between them, the five
emerging-market offices granted 84\% of the invention patents filed in 2000--2005 at the national offices of
emerging economies in PATSTAT, and the five advanced-economy offices 83\% of those filed at the national offices of
advanced economies.\footnote{Advanced economies are those the World Bank classed as high income throughout 2000--2005
\citep{WBincome}, including Taiwan (high income since 1987); regional offices such as the European Patent Office are
left out. Korea, classed as high income in 1995--97 and again from 2001, is counted as emerging, since it moved
between the two groups during the sample period. Without it, small firms are level ($-0.2$ points), mid-sized and large firms are 5.8 and 4.8 points
ahead, and universities are 18.2 points behind.}

\paragraph{Owners.} An owner is domestic if the first applicant is from the office's country. Missing applicant
countries (common in China and Japan) are imputed from the patent family and the office of first filing. Owners are
grouped by sector and size. Size is the number of distinct inventions (DOCDB families) the owner
filed worldwide in 2000--2005.\footnote{An owner combines harmonised applicants that share a name, a distinctive
name prefix or most of their patent families, so that local subsidiaries of multinationals are sized with their parent (Appendix~\ref{app:data}).} Small firms have 1--9
inventions, mid-sized firms 10--99 and large firms 100 or more; universities and public research organisations form a
fourth group. Individual inventors are excluded: their patents behave differently, and in China they were the main
target of local filing subsidies in this period. So are patents whose owner cannot be classed as a firm or a
university. Also, patents granted more than ten years after filing are dropped, except in Brazil, where the median
grant lag was 11 years; there every patent whose survival eight years after grant is observed by mid-2025 is kept.
Table~\ref{tab:sample} describes the sample: patents by office and owner in panel A, owners in panel B. It covers
162{,}640 small firms, 36{,}815 mid-sized firms, 23{,}624 large firms and 17{,}060 universities and public research
organisations, each owner counted once whether it is domestic or foreign.

\begin{table}[t]
\centering
\begin{threeparttable}
\caption{The sample: granted patents filed 2000--2005, firms and universities}\label{tab:sample}
\small
% Generated by em-patent-value code/38_paper_exhibits.py (sample_counts.csv, owner_counts.csv); do not edit by hand
\setlength{\tabcolsep}{4pt}
\begin{tabular}{lrrrrrrrrr}
\toprule
 & & \multicolumn{4}{c}{Domestic owners} & \multicolumn{4}{c}{Foreign owners} \\
\cmidrule(lr){3-6}\cmidrule(lr){7-10}
 & & \multicolumn{3}{c}{Firms} & & \multicolumn{3}{c}{Firms} & \\
\cmidrule(lr){3-5}\cmidrule(lr){7-9}
Office & Patents & Small & Mid & Large & Univ. & Small & Mid & Large & Univ. \\
\midrule
\multicolumn{10}{l}{\emph{A. Patents, by office}} \\
China          & 336{,}371 & 24{,}621 & 10{,}804 & 31{,}124 & 39{,}836 & 23{,}354 & 19{,}765 & 182{,}885 & 3{,}982 \\
Korea          & 308{,}140 & 28{,}157 & 22{,}758 & 161{,}619 & 19{,}530 & 3{,}381 & 6{,}571 & 64{,}876 & 1{,}248 \\
Russia         & 49{,}912 & 9{,}079 & 4{,}264 & 1{,}025 & 15{,}738 & 6{,}114 & 1{,}116 & 12{,}109 & 467 \\
Taiwan         & 149{,}886 & 5{,}753 & 7{,}927 & 28{,}968 & 6{,}665 & 4{,}595 & 9{,}828 & 85{,}234 & 916 \\
Czech Rep.     & 3{,}674 & 568 & 340 & 1 & 95 & 436 & 708 & 1{,}502 & 24 \\
\addlinespace
United States  & 938{,}481 & 67{,}629 & 80{,}826 & 291{,}287 & 28{,}373 & 38{,}183 & 51{,}046 & 364{,}135 & 17{,}002 \\
Japan          & 677{,}557 & 20{,}878 & 49{,}397 & 548{,}159 & 15{,}540 & 2{,}293 & 4{,}053 & 36{,}096 & 1{,}141 \\
Germany        & 81{,}819 & 10{,}209 & 13{,}751 & 35{,}886 & 3{,}587 & 1{,}360 & 1{,}868 & 14{,}948 & 210 \\
United Kingdom & 41{,}157 & 4{,}954 & 4{,}164 & 2{,}792 & 445 & 2{,}571 & 4{,}255 & 21{,}385 & 591 \\
France         & 58{,}950 & 9{,}722 & 9{,}858 & 23{,}665 & 3{,}753 & 1{,}187 & 1{,}460 & 9{,}143 & 162 \\
\midrule
Total          & 2{,}645{,}947 & 181{,}570 & 204{,}089 & 1{,}124{,}526 & 133{,}562 & 83{,}474 & 100{,}670 & 792{,}313 & 25{,}743 \\
\midrule
\multicolumn{10}{l}{\emph{B. Owners, by office}} \\
China          &   & 18{,}126 & 1{,}235 & 648 & 6{,}122 & 20{,}494 & 8{,}919 & 10{,}609 & 1{,}738 \\
Korea          &   & 16{,}782 & 3{,}441 & 814 & 732 & 2{,}788 & 3{,}355 & 4{,}786 & 433 \\
Russia         &   & 5{,}780 & 487 & 35 & 3{,}679 & 4{,}181 & 199 & 855 & 289 \\
Taiwan         &   & 3{,}423 & 616 & 119 & 183 & 3{,}576 & 3{,}200 & 2{,}983 & 270 \\
Czech Rep.     &   & 368 & 23 & 1 & 55 & 378 & 442 & 408 & 23 \\
\addlinespace
United States  &   & 38{,}861 & 7{,}298 & 2{,}841 & 1{,}369 & 27{,}864 & 9{,}766 & 5{,}121 & 2{,}129 \\
Japan          &   & 13{,}986 & 6{,}635 & 3{,}261 & 497 & 1{,}760 & 1{,}750 & 1{,}958 & 282 \\
Germany        &   & 6{,}813 & 2{,}106 & 643 & 341 & 1{,}052 & 606 & 963 & 79 \\
United Kingdom &   & 3{,}497 & 804 & 190 & 145 & 2{,}015 & 1{,}448 & 1{,}154 & 210 \\
France         &   & 5{,}859 & 993 & 278 & 296 & 919 & 563 & 697 & 60 \\
\midrule
All offices    &   & 113{,}495 & 23{,}630 & 8{,}783 & 13{,}410 & 59{,}249 & 20{,}318 & 18{,}610 & 4{,}510 \\
\bottomrule
\end{tabular}
\begin{tablenotes}\footnotesize
\item \emph{Notes:} Granted invention patents filed 2000--2005 with a grant lag of up to ten years; individual owners and owners of unknown type excluded. Firm size: inventions (DOCDB families) filed worldwide in 2000--2005 by the owner's corporate group (harmonised applicants linked by name and by shared patent families); small 1--9, mid-sized 10--99, large 100 or more. Univ.: universities and public research organisations. Panel B: distinct harmonised applicants with at least one sample patent in the office, classed by the size of their group. An owner patenting in several offices is counted in each; All offices counts it once among domestic and once among foreign owners, so the two totals overlap and fall short of the sums over offices. Source: PATSTAT Global, autumn 2025.
\end{tablenotes}
\end{threeparttable}
\end{table}




\paragraph{Quality.} Observable quality is measured by forward citations to the patent family, in seven bins (0, 1, 2,
3--4, 5--9, 10--19, 20 or more),\footnote{Bins of roughly doubling width, so that domestic patents
are compared with foreign patents of similar citations rather than through a fitted slope. Office-specific linear
slopes in log citations give almost the same estimates (Appendix Table~\ref{tab:rich}).} and the WIPO technology field (35 fields). Domestic patents in emerging-market offices are cited far less than foreign
ones, by a factor of five in China and six in Korea, so quality controls matter.

\paragraph{Design.} The basic comparison is between domestic and foreign owners in the same office. Foreign owners are likely multinationals protecting a foreign market, so domestic
and foreign patents differ everywhere; the question is whether the domestic shortfall is larger in emerging-market
offices. For patent $i$ of owner group $g$ in office $o$, filed in year $t$ in field $f$, in citation bin $c$ and granted
after $\ell_i$ years,
\begin{equation}\label{eq:did}
  y_{i} = \beta\, D_i + \delta\, D_i \times EM_o + \zeta\, D_i (\ell_i - \bar\ell) + \mu_{ot} + \lambda_{of}
  + \kappa_{og} + \psi_{oc} + \nu_{o\lfloor\ell_i\rfloor} + \varepsilon_i ,
\end{equation}
where $y_i$ indicates that the patent is alive eight years after grant, $D_i$ marks a domestic owner and $EM_o$ an
emerging-market office; $\bar\ell$ is the mean grant lag, so that $\beta$ is the domestic--foreign gap in the
advanced-economy offices at the mean grant lag and $\zeta$ lets that gap vary with the grant lag. The remaining terms
are office-specific effects of filing year ($\mu_{ot}$), technology field ($\lambda_{of}$), owner group
($\kappa_{og}$, in the pooled regression only), citation bin ($\psi_{oc}$) and grant lag in whole years
($\nu_{o\lfloor\ell_i\rfloor}$); they absorb the main effect of $EM_o$.\footnote{Survival is measured from grant because renewal fees are due from grant. Foreign
owners' patents are granted 2--4 years later than domestic ones in the emerging-market offices (Appendix
Figure~\ref{fig:lags}), so at a given time after grant they are older, face higher fees where fees rise with age from filing, and
have less of their term left. Comparing patents with the same grant lag removes this advantage of domestic owners.}
The coefficient of interest, $\delta$, is the extra domestic--foreign gap in emerging-market offices, in percentage
points. Standard errors are clustered by owner. Because $\delta$ compares five offices with five, Appendix Table~\ref{tab:rich} also
reports inference across offices.
